\section*{Chapter \ref{ch:g12:synthesis-strategy} --- Synthesis Strategy and Green Chemistry}

\begin{solution}{exo:g12:synthesis-strategy:1}
\omterm{def:g11:dissolution:solvation}{Hydration}: $46.0/(28.0 + 18.0) = \qty{100}{\percent}$. Fermentation:
$2 \times 46.0/180.0 = \qty{51.1}{\percent}$.
\end{solution}

\begin{solution}{exo:g12:synthesis-strategy:2}
Both groups, \ce{Z} on the \omterm{def:g11:functional-groups:amine}{amine} of glycine and the methyl \omterm{def:g11:functional-groups:carboxylic-acid}{ester} on the acid
of alanine, are put on at step 1 and taken off at step 3.
\end{solution}

\begin{solution}{exo:g12:synthesis-strategy:3}
Principle 5: use safer \omterm{def:g6:solutions-and-solubility:solute}{solvents} and reaction conditions.
\end{solution}

\begin{solution}{exo:g12:synthesis-strategy:4}
An excess of ethanol makes $Q < K$, so the reaction goes further
forwards. A \omterm{def:g12:catalysis:catalyst}{catalyst} speeds up both directions alike: the same \omterm{def:g10:reaction-progress-table:system}{final
state} is reached, only sooner.
\end{solution}

\begin{solution}{exo:g12:synthesis-strategy:5}
Yes: only the \omterm{def:g11:functional-groups:carbonyl}{ketone} is reduced, into a secondary \omterm{def:g11:functional-groups:alcohol}{alcohol} group; the
\omterm{def:g11:functional-groups:carboxylic-acid}{ester} is unchanged.
\end{solution}

\begin{solution}{exo:g12:synthesis-strategy:6}
More \omterm{def:g11:functional-groups:carboxylic-acid}{ester}: an excess of one \omterm{def:g7:chemical-reactions:reactant}{reactant}, or removing the water (water trap)
as it forms. Faster: heat under reflux, add an acid \omterm{def:g12:catalysis:catalyst}{catalyst}.
\end{solution}

\begin{solution}{exo:g12:synthesis-strategy:7}
With \ce{NaBH4}:
\[ \ce{CH3-CH(OH)-CH2-CH2-CO-O-CH3} . \]
With \ce{LiAlH4}:
\[ \ce{CH3-CH(OH)-CH2-CH2-CH2OH} \]
(pentane-1,4-diol), plus methanol from the \omterm{def:g11:functional-groups:carboxylic-acid}{ester}'s other half.
\end{solution}

\begin{solution}{exo:g12:synthesis-strategy:8}
The addition (\qty{100}{\percent}), the esterification
(\qty{83}{\percent}) and the three-step route to ibuprofen
(\qty{77}{\percent}). $77.4/40.0 = 1.9$: almost twice as good.
\end{solution}

\begin{solution}{exo:g12:synthesis-strategy:9}
$180.0/(138.0 + 102.0) = \qty{75.0}{\percent}$. Ethanoic acid:
$60.0/180.0 = \qty{0.333}{kg}$ per kilogram of aspirin.
\end{solution}

\begin{solution}{exo:g12:synthesis-strategy:10}
The heating speeds up the reaction; the condenser returns the vapours
to the flask, so that no \omterm{def:g7:chemical-reactions:reactant}{reactant} or \omterm{def:g6:solutions-and-solubility:solute}{solvent} is lost. Heating improves
the rate, not the final \omterm{def:g10:synthesis-yield:yield}{yield}.
\end{solution}

\begin{solution}{exo:g12:synthesis-strategy:11}
Six steps against three. Principles: prevent waste (1), maximise \omterm{def:g12:synthesis-strategy:atom-economy}{atom
economy} (2), use \omterm{def:g12:catalysis:catalyst}{catalysts} (9); also fewer steps, less energy (6).
\end{solution}

\begin{solution}{exo:g12:synthesis-strategy:12}
$0.80 \times 0.70 \times 0.90 = 0.504$: \qty{50.4}{\percent}. Starting
material: $0.10/0.504 = \qty{0.20}{mol}$.
\end{solution}

\begin{solution}{exo:g12:synthesis-strategy:13}
Protect the \omterm{def:g11:functional-groups:amine}{amine} of alanine (\ce{Z}) and the acid of glycine (methyl \omterm{def:g11:functional-groups:carboxylic-acid}{ester});
form the \omterm{def:g11:functional-groups:amine}{amide} bond between the free acid of alanine and the free \omterm{def:g11:functional-groups:amine}{amine}
of glycine; remove both \omterm{def:g12:synthesis-strategy:protecting-group}{protecting groups}. Three stages, that is five
reactions if each protection and deprotection is counted.
\end{solution}

\begin{solution}{exo:g12:synthesis-strategy:14}
The \omterm{def:g12:catalysis:catalyst}{catalyst} applies principle 9. But \qty{250}{\celsius} goes against
principle 6, benzene against principles 3 and 5, and an \omterm{def:g12:synthesis-strategy:atom-economy}{atom economy} of
\qty{35}{\percent} against principle 2: $65/35 = \qty{1.9}{kg}$ of waste
per kilogram of product. One green feature does not make a green
process.
\end{solution}

\begin{solution}{exo:g12:synthesis-strategy:15}
% ledger: crit:CO2-T, crit:CO2-p
\qty{40}{\celsius} and \qty{100}{bar} are above the critical point of
carbon dioxide, \qty{31}{\celsius} and \qty{73.8}{bar}. The pressure is a
hundred times atmospheric pressure: the vessel is thick steel. Carbon
dioxide is neither toxic nor flammable, and it turns back into a gas
leaving no residue, and can be reused; no chlorinated \omterm{def:g6:solutions-and-solubility:solute}{solvent} escapes.
\end{solution}

\begin{solution}{pb:g12:synthesis-strategy:1}
% ledger: ibu:bhc-ae-recovered
\textbf{1.} C: $10 + 4 + 1 = 15 = 13 + 2$; H: $14 + 6 + 2 = 22 = 18 + 4$;
O: $3 + 1 = 4 = 2 + 2$.

\textbf{2.} Ethanoic acid, \ce{C2H4O2}.

\textbf{3.} Nitrogen, chlorine and sodium: ibuprofen holds only carbon,
hydrogen and oxygen.

\textbf{4.} The newer one: principle 9, \omterm{def:g12:catalysis:catalyst}{catalysts} rather than
stoichiometric \omterm{def:g7:identifying-substances:characteristic-test}{reagents}.

\textbf{5.} $13 \times 12.0 + 18 \times 1.0 + 2 \times 16.0 =
\qty{206.0}{g/mol}$.

\textbf{6.} $134.0 + 102.0 + 122.5 + 68.0 + 19.0 + 33.0 + 36.0 =
\qty{514.5}{g/mol}$.

\textbf{7.} $206.0/514.5 = \qty{40.0}{\percent}$.

\textbf{8.} $134.0 + 102.0 + 2.0 + 28.0 = \qty{266.0}{g/mol}$.

\textbf{9.} $206.0/266.0 = \qty{77.4}{\percent}$.

\textbf{10.} $(206.0 + 60.0)/266.0 = \qty{100}{\percent}$ on paper; over
\qty{99}{\percent} in practice.

\textbf{11.} $514.5/206.0 = \qty{2.50}{t}$.

\textbf{12.} $2.50 - 1 = \qty{1.50}{t}$.

\textbf{13.} $266.0/206.0 = \qty{1.29}{t}$.

\textbf{14.} $60.0/206.0 = \qty{0.29}{t}$.

\textbf{15.} $1.50 - 0.29 = \qty{1.21}{t}$.

\textbf{16.} $0.90^6 = 0.53$: \qty{53}{\percent}.

\textbf{17.} $0.90^3 = 0.73$: \qty{73}{\percent}.

\textbf{18.} Each step uses \omterm{def:g6:solutions-and-solubility:solute}{solvents}, separations and purifications, and
loses some product: fewer steps, fewer of these losses.

\textbf{19.} Almost none: its only by-product is used.

\textbf{20.} The three-step route avoids about \qty{1.5}{t} of waste per
tonne of ibuprofen (\qty{1.2}{t} even if its ethanoic acid were thrown
away).
\end{solution}
